%  \indent Układem mechanicznym wykorzystanym w pracy jest aluminiowa konstrukcja tzw. goniometru, stanowiącego precyzyjny uchwyt o dwóch stopniach swobody, którego konfigurację można ustawiać aktywnie za pomocą silników. Uchwyt ten z założenia dostosowany jest do instrumentu optycznego i pozwala na korygowanie orientacji jego osi optycznej wyrażonej za pomocą dwóch kątów Eulera. Konstrukcja została zaprojektowana i wykonana w 2021 r. m.in. w ramach rektorskiego projektu międzywydziałowego SiGR. Wykonany wtedy prototypowy układ sterowania, wykorzystujący elementy automatyki przemysłowej firmy B\&R pozwolił na wstępną ocenę właściwości mechanicznych urządzenia. 

% W tej pracy inżynierskiej opracowano i wykonano prototyp nowego układu sterowania, który docelowo pozwoli ograniczyć koszty, zmniejszyć gabaryty i zapewnić lepszą funkcjonalność niż rozwiązanie oparte na sterownikach przesyłowych. W tym celu zaproponowano użycie zaawansowanego 32-bitowego mikrokontrolera, który odpowiada za obsługę głowic pomiarowych, oblicza algorytm sterujący napędami oraz zapewnia interfejs komunikacyjny za pomocą sieci Ethernet. Opracowany układ sterowania w zamiarze jest więc elementem większego systemu, który może nadzorować pracę goniometru zarówno przez operatora jak i w systemie automatycznym.

% %Z racji na całkiem spory zakres obowiązków pracę podzielono na dwie części. 
% W ramach pracy wyodrębniono dwa główne obszary, które strukturyzują zakres obowiązków związane z projektowaniem sprzętu i oprogramowania. Pierwszy z nich obejmował pracę nad odpowiednim połączeniem i doborem elementów wykonawczych oraz zasilających. Następnie przeprowadzono wstępne testy wykorzystywanych urządzeń i sprawdzono zasadę ich działania na podstawie dokumentacji technicznych. Zaprojektowano także schematy połączeń elektrycznych. W międzyczasie rozwijano również drugą część, w której tworzono aplikację webową jako wspomniany wcześniej interface użytkownika do zadawania pozycji docelowej i możliwości doglądania wyników od sprzężenia zwrotnego. W kolejnych etapach z wykorzystaniem odpowiednich protokołów komunikacyjnych nawiązano połączenia z urządzeniami zamontowanymi na stanowisku oraz napisano cały kod programu obsługiwanego przez 
% mikrokontroler. 

% \newpage 
% \indent Aleksander Hanczewski określił %zidentyfikował 
% wymagania związane ze sterowaniem silników i opracował rozwiązanie zapewniające poprawny odczyt pozycji z użyciem enkoderów. Na tej podstawie zaprojektował algorytm sterowania umożliwiający precyzyjne utrzymanie pozycji zadanej. Następnie napisał kod odpowiedzialny za realizację tego algorytmu. Dodatkowo, przygotował schemat połączeń elektrycznych oraz schemat sterowania, okablował urządzenia, a także zajął się montażem pozostałych elementów stanowiska. \newline
% \indent Dominik Gałczyk zaprojektował i stworzył aplikację webową, która obsługuje stanowisko badawcze. Wykorzystując system czasu rzeczywistego FreeRTOS przygotował serwer na mikrokontrolerze. Opracował oraz zaimplementował mechanizm wymiany wiadomości do komunikacji między serwerem, WebSocketem oraz aplikacją. Zrealizował WebSocket pełniący rolę pośrednika między interfejsem użytkownika, a serwerem na mikrokontrolerze. Był także odpowiedzialny za integrację pozostałych modułów kodu. Dodatkowo skonteneryzował aplikację webową przy użyciu technologii Docker.

% \newpage